\chapter{Operating Tellervo Server}
\label{txt:runningserver}

Tellervo Server normally needs little day-to-day attention, but it must be
backed up, patched, monitored, and tested like any other database service. The
\verb|tellervo-server| command is the supported interface for configuring the
webservice and maintaining Tellervo databases. Run administrative commands with
\verb|sudo|.

\section{Checking an installation}

Run the configuration tests after installation, upgrades, certificate changes,
or database maintenance:

\code{sudo tellervo-server --test}

For a named instance:

\code{sudo tellervo-server --instance lab-a --test}

The test checks the installed components, configuration files, Apache setup,
database connection, and webservice. Public reachability may be reported as a
warning because DNS, HTTPS, firewalls, and reverse proxies can be managed
outside the Tellervo package.

Useful diagnostic commands include:

\begin{description}
 \item[\texttt{--version}] Show the installed webservice and database version.
 \item[\texttt{--info}] Show configuration information for the selected
 instance.
 \item[\texttt{--network}] Check and, where possible, repair network settings.
 \item[\texttt{--repair-web-assets}] Restore packaged webservice files without
 replacing instance-specific configuration.
 \item[\texttt{--sysconfig}] Regenerate \verb|systemconfig.php| from the
 instance configuration.
 \item[\texttt{--emptylog}] Clean Tellervo log files and database log tables.
\end{description}

Use \verb|sudo tellervo-server --help| as the authoritative list for the
installed release.

\section{Backing up and restoring}
\index{Server!Backup}
\index{Backup}

A complete Tellervo backup has at least two parts:

\begin{enumerate}
 \item A PostgreSQL database dump for every Tellervo instance.
 \item A file-level backup of uploaded ODK media, linked files, instance
 configuration, and any locally managed TLS or reverse-proxy configuration.
\end{enumerate}

Store backups on a different machine or managed backup service. A backup kept
only on the Tellervo host will be lost with that host. Encrypt backups that
contain sensitive research, location, or user information, restrict access to
them, and test restoration regularly.

\subsection{Database backup}
\label{txt:dbbackup}

Back up the default instance to an explicitly named destination:

\code{sudo tellervo-server --backup-db /srv/backups/tellervo.dump}

For a named instance:

\code{sudo tellervo-server --instance lab-a --backup-db /srv/backups/lab-a.dump}

Use a scheduler such as a systemd timer, cron, or an institutional backup agent
to run backups automatically. Include the date and instance in each filename,
retain multiple generations, monitor failures, and periodically copy a backup
to a restore-test host. Do not use \verb|/tmp| for durable backups.

\subsection{Media and configuration backup}

Database dumps do not contain uploaded media files. Back up each configured
instance web folder and the instance registry:

\begin{itemize}
 \item \verb|/etc/tellervo-server/instances.conf|
 \item the web folder shown by \verb|tellervo-server --instance NAME --info|
 \item locally managed Apache, proxy, firewall, and TLS configuration
\end{itemize}

Credentials and private keys require stricter permissions than ordinary media
files. Follow the backup policy of the institution rather than copying secrets
to an unmanaged synchronization account.

\subsection{Restoring a database}
\label{txt:dbrestore}

\warn{A restore replaces the selected instance database. Confirm the instance
name and backup file before continuing, and keep a backup of the current
database until the restored service has been verified.}

Restore the default instance with:

\code{sudo tellervo-server --restore-db /srv/backups/tellervo.dump}

Restore a named instance with:

\code{sudo tellervo-server --instance lab-a --restore-db /srv/backups/lab-a.dump}

After restoration, replace the matching media files if necessary, apply any
outstanding database upgrades, regenerate the system configuration, and test:

\code{sudo tellervo-server --instance lab-a --upgrade-db}
\code{sudo tellervo-server --instance lab-a --sysconfig}
\code{sudo tellervo-server --instance lab-a --test}

\section{Upgrading the server}
\label{txt:upgradeserver}
\index{Upgrading!Tellervo server}

Before every upgrade:

\begin{enumerate}
 \item Read the release notes for operating-system, PostgreSQL, and desktop
 compatibility changes.
 \item Back up every database and the corresponding media files.
 \item Confirm that a recent backup can be read on another machine.
 \item Plan a maintenance window and stop desktop users from editing data.
\end{enumerate}

For a server installed from an APT repository:

\code{sudo apt update}
\code{sudo apt upgrade}
\code{sudo tellervo-server --test}

For local packages, supply the updated common, webservice, database selector,
provider, and meta packages together with:

\code{sudo apt install ./package.deb}

APT resolves dependencies; raw \verb|dpkg --install| does not.

Package hooks normally apply database upgrades and regenerate configuration.
If maintenance must be run explicitly:

\code{sudo tellervo-server --upgrade-db}
\code{sudo tellervo-server --sysconfig}
\code{sudo tellervo-server --test}

For several registered instances, run the corresponding cluster operation:

\code{sudo tellervo-server --cluster upgrade-db}
\code{sudo tellervo-server --cluster sysconfig}
\code{sudo tellervo-server --cluster test}

Do not perform a major PostgreSQL or operating-system upgrade as if it were an
ordinary Tellervo package update. Follow the operating-system and PostgreSQL
upgrade guidance, retain the old cluster until validation is complete, and test
Tellervo against the upgraded database before reopening it to users.

\section{Named instances and clusters}

A single packaged web host can serve several independent Tellervo databases.
Each named instance has its own database name, web folder, credentials file, and
Apache URL alias. This is useful when several laboratories share one managed
server but must keep their data separate.

Create or configure an instance:

\code{sudo tellervo-server --instance lab-a --configure}

Unless overridden, the instance uses a database such as
\verb|tellervo_lab_a|, a web folder such as \verb|/var/www/lab-a/|, and URL path
\verb|/lab-a/|. Explicit values can be supplied when local naming policy
requires them:

\begin{lstlisting}
sudo tellervo-server \
  --instance lab-a \
  --dbname tellervo_lab_a \
  --webfolder /var/www/lab-a \
  --configure
\end{lstlisting}

Configured instances are registered in
\verb|/etc/tellervo-server/instances.conf|. Common commands are:

\code{sudo tellervo-server --list-instances}
\code{sudo tellervo-server --cluster test}
\code{sudo tellervo-server --cluster version}
\code{sudo tellervo-server --cluster info}

Cluster operations support \verb|test|, \verb|sysconfig|,
\verb|upgrade-db|, \verb|version|, \verb|info|, \verb|emptylog|, and
\verb|network|.

\warn{\texttt{tellervo-server --delete-instance NAME} drops the named database
and removes its web folder. It is intentionally interactive. Make and verify
database and media backups before using it.}

\section{Split web and database hosts}

In a split installation the database host runs the common and database
packages, while the application host runs the common and webservice packages.
Section \ref{txt:splitinstall} gives the complete initial setup procedure,
including creation of the role and database, loading the packaged template,
and testing the connection from the web host.

PostgreSQL should listen only on the required interfaces. Add a
\verb|pg_hba.conf| rule limited to the web host, database, and role. For
example:

\begin{lstlisting}
host  tellervo_lab_a  tellervo_lab_a  192.0.2.20/32  scram-sha-256
\end{lstlisting}

An HBA-only change can be reloaded without interrupting the server:

\code{sudo systemctl reload postgresql}

Changing \verb|listen_addresses| requires a restart instead. After any network
or authentication change, repeat the direct \verb|psql| connection test from
the web host before troubleshooting Tellervo itself. Then rerun the instance
tests:

\code{sudo tellervo-server --instance lab-a --test}

The remote wizard writes the connection configuration and tests the existing
database; it does not create a database remotely. Permit TCP port 5432 only
between the necessary hosts, and use a protected network or PostgreSQL TLS when
traffic crosses an untrusted network.

\section{Security}
\index{Security}

\begin{itemize}
 \item Keep the operating system, Apache, PHP, PostgreSQL, and Tellervo packages
 patched.
 \item Publish the webservice through HTTPS with a currently valid certificate.
 Do not expose a production Tellervo login over plain HTTP.
 \item Restrict PostgreSQL to the webservice host and administrator networks.
 PostgreSQL should not be generally reachable from the Internet.
 \item Use individual Tellervo accounts, disable accounts that are no longer
 needed, and grant administrative group membership sparingly.
 \item Use strong, unique passwords for system, database, and Tellervo accounts.
 Never retain installer or demonstration credentials.
 \item Protect configuration files and backups because they contain credentials
 or sensitive data.
 \item Review Apache, PHP, PostgreSQL, Tellervo, and backup logs and monitor disk
 capacity.
\end{itemize}

Tellervo's legacy challenge-response protocol includes MD5-based password
hashing. HTTPS is therefore a required outer security boundary for production
deployments, not an optional enhancement. A reverse proxy or institutional load
balancer may terminate HTTPS, but traffic between that proxy and the
webservice must also be confined to a trusted network.

\section{Direct database access}
\index{Database}
\index{PostgreSQL}

Tellervo is designed to access PostgreSQL through its webservice. Direct
read-only SQL can be useful for reporting, but application tables and functions
form an internal interface that can change during upgrades.

\warn{Do not insert, update, or delete Tellervo records directly. Bypassing the
webservice can violate hierarchy, permission, versioning, and audit assumptions
and may make later database upgrades fail. Take a backup before investigating a
production database.}

Use a current PostgreSQL client such as \verb|psql| or pgAdmin. Do not weaken
\verb|pg_hba.conf| with a broad ``all users/all databases'' rule. Create a
read-only reporting role where possible, restrict its source addresses, and
prefer SSH tunnelling or TLS rather than exposing PostgreSQL publicly.

\section{Server configuration}
\label{txt:serverConfig}

Run the interactive configuration wizard for the default instance with:

\code{sudo tellervo-server --configure}

Use \verb|--reconfigure| to overwrite an existing generated configuration.
Select a named instance by adding \verb|--instance NAME|. Other maintenance
options include:

\begin{itemize}
 \item \verb|--set-db-pwd| to change the database role password and regenerate
 the webservice configuration
 \item \verb|--set-admin-pwd| to reset the Tellervo administrator password
 \item \verb|--start|, \verb|--stop|, and \verb|--restart| for packaged services
 \item \verb|--cascade-delete true| or \verb|false| to control the configured
 delete policy
\end{itemize}

Avoid editing generated \verb|systemconfig.php| files directly; regenerate them
with the administration command so that instance configuration remains
consistent.

\section{Managing map services}
\label{txt:managingmaps}
\index{Web Map Service (WMS)}

Add a shared Web Map Service to the selected Tellervo database with:

\begin{lstlisting}
sudo tellervo-server \
  --add-wms-name "Service name" \
  --add-wms-url "https://maps.example.org/wms"
\end{lstlisting}

Remove it with:

\code{sudo tellervo-server --del-wms-name "Service name"}

Add \verb|--instance NAME| when managing a named instance. Test third-party map
services before advertising them to users; availability, authentication,
supported layers, and usage terms are controlled by the map provider.
